% Initial setting
\documentclass{article}
\usepackage{graphicx} 
\usepackage{kotex}
\title{제1장 시계열의 정상성 및 시계열 분석의 기초}
\author{Minjae Kim}
\date{} % date 비워놓으면 날짜 안 뜸
\renewcommand{\thesection}{\arabic{section}.}
\renewcommand{\thesubsection}{\arabic{section}.\arabic{subsection}.}
\renewcommand{\thesubsubsection}{\arabic{section}.\arabic{subsection}.\arabic{subsubsection}.}
\usepackage{setspace}
\usepackage{booktabs}
\usepackage{float}
\usepackage{adjustbox}
\usepackage[table]{xcolor}
\usepackage{longtable}
\usepackage{makecell}
\usepackage{pdflscape}
\usepackage{placeins}
\usepackage{caption}
\usepackage{silence}
\WarningFilter{*}{*}
\usepackage[a4paper, margin=2cm]{geometry}
\usepackage{amsmath}
\usepackage{amsthm}
\newtheorem{defn}{Definition}
\usepackage{amssymb}
\hbadness=10000
\begin{document}
\sloppy
\doublespacing
\maketitle
% 1. 시계열의 정상성
\section{시계열의 정상성}
  \begin{itemize}
    \item 시계열 자료란?
      \begin{itemize}
        \item 동일한 개체를 여러 시간에 걸쳐 관측하여 수집한 자료를 의미함.
        \item 엄밀하게 시계열을 이해하기 위해서는 확률과정(stochasitc process)의 개념이 필요하다.\footnote{이에 대한 내용은 추후 계량경제학, 시계열 분석 등을 학습할 때 정리해보도록 하자.} 
      \end{itemize}
  \end{itemize}
\begin{defn}[정상성]\label{def:stationarity}
\emph{정상시계열}(stationary time series)은 시계열 데이터의 특성이 시간에 따라 일정하게 유지되는 성질을 의미한다. 
\end{defn}
  \begin{itemize}
    \item 정상성은 강정상성(strict stationarity)과 약(또는 공분산)정상성(weak or covariance stationarity)으로 구분되는데, 강정상성을 보이기는 다소 어려운 과정이므로 실증분석에서는 공분산 정상성을 만족함을 요구하는 것이 일반적이다.
    \item 아래와 같이, 강정상성과 약정상성의 정의를 정리해보자.
  \end{itemize}
\begin{defn}[강정상성]\label{def:strict stationarity}
임의의 $s$에 대하여 ($y_{t_{1}}$, \dots , $y_{t_{m}}$)과 ($y_{t_{1} + s}$, \dots , $y_{t_{m} + s}$)의 결합확률분포가 동일할 때,
시계열 \{$y_{t}$\}는 \emph{강정상성을 만족하는 시계열}(strictly stationary time series)이다.
\end{defn}
\begin{defn}[약 또는 공분산 정상성]\label{def:weak(covariance) stationarity}
시계열 \{$y_{t}$\}는 다음 세 가지 조건이 만족한다면 \emph{공분산 정상시계열}(covariance stationary time series)이라 부른다.
  \begin{enumerate}
    \item 시계열의 평균이 시간에 따라 일정하다. $E(y_t) = E(y_{t+s}), \ \forall s$
    \item 시계열의 분산이 시간에 따라 일정하다. $var(y_t) = var(y_{t+s}), \ \forall s$
    \item 두 시점 $t$와 $t-k$의 공분산 $cov(y_{t}, y_{t-k})$이 시간 $t$에는 의존하지 않고, 시차 $k$만의 함수로 표현된다. $cov(y_t, y_{t-k}) = cov(y_{t+s}, y_{t+s-k}), \ \forall s, k$
  \end{enumerate}  
\end{defn}
시계열 예측에서 대부분은 과거 자료의 특성을 학습한 후 동일한 특성이 유지될 것으로 가정하여 미래를 예측한다.\\
따라서 시계열 예측에서 정상성은 주의해야 할 매우 중요한 개념이다.
\newpage
% 2. 기본 시계열 모형
\section{기본 시계열 모형}
%% 2.1. 백색잡음, 자기회귀, 이동평균
\subsection{백색잡음, 자기회귀, 이동평균}
\begin{defn}[백색잡음]\label{def:white noise}
분산이 일정하고 자기상관이 없는 확률과정을 \emph{백색잡음}(white noise)이라 한다.  
\end{defn}
  \begin{itemize}
    \item 평균이 0이고 분산이 $\sigma^{2}$인 백색잡음을 $\{\varepsilon_{t}\} \sim WN(0, \sigma^2)$으로 나타낸다.
    \item 이 경우 $\{\varepsilon_t\}$는 정상시계열임을 쉽게 확인할 수 있다.\footnote{평균과 분산이 일정한 상수이고, 자기공분산도 $k$가 0일 때와 아닐 때 $\sigma^2$과 0으로 $k$에 의존하기 때문이다.}
  \end{itemize}
\begin{defn}[1차 자기회귀과정]\label{def:AR1}
  시계열 $y_{t}$는 다음 세 가지 조건이 만족할 때, 1차 자기회귀과정(autoregressive process)이라 부른다.\footnote{AR(1)과정으로 간단하게 표현할 수 있다.}
    \begin{enumerate}
      \item $y_{t} = \alpha + \beta y_{t-1} + \varepsilon_{t}$
      \item $\{\varepsilon_t\} \sim WN(0, \sigma^2)$
      \item $\varepsilon_t$는 $y_t$의 모든 과거값과 상관관계가 없다.
    \end{enumerate}
\end{defn}
위와 같이 정의된 AR(1)과정인 $\{y_t\}$의 통계적 특성들을 다음과 같이 도출해보자.\footnote{$\{y_t\}$는 안정적인 시계열이라 가정한다.}
\subsubsection*{$\checkmark$ $E(y_t) = \frac{\alpha}{1-\beta} \ or \ 0$ }
먼저, 상수항을 포함하는 일반적인 AR(1)과정 $\{y_t\}$의 expectation을 유도해보자.\\
$y_{t} = \alpha + \beta y_{t-1} + \varepsilon_t$, $\{\varepsilon_t\} \sim WN(0, \sigma^2)$에서 안정적인 시계열의 특성을 이용하면, $E(y_t) = E(y_{t-1})\equiv \mu$. \\
\begin{equation}
\begin{aligned}
E(y_t) &= \alpha + \beta E(y_{t-1}) + E(\varepsilon_t) \\
       &= \alpha + \beta E(y_{t-1}) \ (\because E(\varepsilon_t) = 0)
\end{aligned}
\end{equation}
식 (1)에서 $E(y_{t})$과 $E(y_{t-1})$은 둘 다 상수인 $\mu$로 설정되었으므로, $\mu = \alpha + \beta\mu$가 된다.\\
따라서, $\mu = \frac{\alpha}{1-\beta}$.\\
이번에는 상수항을 포함하지 않는 AR(1)과정에서 기대값을 유도해보자. 이 때 $y_t - \mu \equiv x_t $ \\
아래의 식 (2)는 상수항을 포함하였던 기존의 AR(1) 과정에서 상수항을 제거하는 과정이다.
\begin{equation}
\begin{aligned}
y_t - \mu &= \alpha + \beta y_{t-1} + \varepsilon_t - \mu \\
    x_{t} &= \alpha + \beta(x_{t-1} + \mu) + \varepsilon_t - \mu \\
          &= \beta x_{t-1} + \varepsilon_t + \alpha + \beta\mu - \mu       
\end{aligned}  
\end{equation}
식 (2)의 $\alpha + \beta\mu - \mu$는 앞서 $E(y_t) = \mu$를 도출하는 과정을 보면 0이 된다는 것을 자연스럽게 파악할 수 있다.\\
$\therefore x_t = \beta x_{t-1} + \varepsilon_t$가 되면서 상수항이 완벽하게 사라짐을 알 수 있다. $y_t - \mu$가 $x_t$인 사실을 통해 기대값이 0이되는 것도 역시 쉽게 이해할 수 있다.
\newpage
\subsubsection*{$\checkmark$ 분산 구하기}
분산의 경우, 상수항 포함 여부와 관계없이 $\frac{\sigma^2}{(1-\beta^2)}$으로 일정함을 알 수 있다.\\
식 (3)과 같은 과정을 통해, $y_{t} = \alpha + \beta y_{t-1} + \varepsilon_t$, $\{\varepsilon_t\} \sim WN(0, \sigma^2)$에서 분산을 구할 수 있다.\\
안정적인 시계열이라고 가정했으므로, $var(y_t) = var(y_{t-1}) \equiv \sigma^2_{y}$라 해보자.
\begin{equation}
\begin{aligned}
var(y_t) &= var(\alpha + \beta y_{t-1} + \varepsilon_t)\\
         &= \beta^2 \sigma^2_{y} + \sigma^2           
\end{aligned}  
\end{equation}
따라서, $var(y_t)(= \sigma^2_{y}) = \frac{\sigma^2}{(1-\beta^2)}$가 됨을 확인할 수 있다.
\subsubsection*{$\checkmark$ 자기공분산 $\gamma(k)$}
자기공분산 $cov(y_t, y_{t-k}) \equiv \gamma(k)$도 상수항 포함 여부에 관계없이 $\beta^{k}\frac{\sigma^2}{(1-\beta^2)}$으로 일정함을 알 수 있다.\\
\textbf{case 1) 상수항이 없는 AR(1) 과정에서의 자기공분산}\\
$x_t = \beta x_{t-1} + \varepsilon_t$에서 $cov(x_{t}, x_{t-k}) \equiv \gamma(k)$로 놓으면,
\begin{equation}
\begin{aligned}
\gamma(k) &= cov(x_{t}, x_{t-k})\\
          &= cov(\beta x_{t-1} + \varepsilon_t, x_{t-k})\\
          &= \beta cov(x_{t-1}, x_{t-k}) + cov(\varepsilon_t, x_{t-k})\\
          &= \beta\gamma(k-1)\\
          &= \\
          \vdots \\
          &= \beta^{k}\gamma(0)     
\end{aligned}  
\end{equation}
위의 식 (4)에서 $\gamma(0) = cov(x_{t-k},x_{t-k}) = var(x_{t-k}) = var(x_t)$가 되어,\\
$\gamma(0) = \frac{\sigma^2}{(1-\beta^{2})}$가 된다. 따라서, $\gamma(k) = \beta^{k}\frac{\sigma^2}{(1-\beta^2)}$임을 알 수 있다.\\
\textbf{case 2) 상수항을 포함한 AR(1) 과정에서의 자기공분산}\\
이번에는 $y_t = \alpha + \beta y_{t-1} + \varepsilon_t$에서의 자기공분산 함수 $\gamma(k)$를 구해보자.\\
\begin{equation}
\begin{aligned}
cov(y_{t}, y_{t-k}) &= cov(\alpha + \beta y_{t-1} + \varepsilon_t, y_{t-k})\\
                    &= cov(\alpha, y_{t-k}) + \beta cov(y_{t-1}, y_{t-k}) + cov(\varepsilon_t, y_{t-k})\\
                    &= \beta\gamma(k-1)\\
                    \vdots \\
                    &= \beta^{k}\gamma(0)\\
                    &= \beta^{k}\frac{\sigma^2}{(1-\beta^2)}   
\end{aligned}
\end{equation}
\newpage
\subsubsection*{$\checkmark$ AR(1) 과정이 안정적인 시계열이 되려면?}
$|\beta|<1$을 만족해야 함. 다음과 같이 $y_{t} = \alpha + \beta y_{t-1} + \varepsilon_t$, $\varepsilon_t \sim WN(0, \sigma^2)$에 대해서 저 조건을 만족하지 못하는 경우를 알아보자.\\
\textbf{case 1) $\beta = 1$, Random walk}\\
\begin{equation}
\begin{aligned}
y_{t} &= \alpha + y_{t-1} + \varepsilon_t\\
      &= \alpha + (\alpha + y_{t-2} + \varepsilon_{t-1}) + \varepsilon_t \\
      &= 2\alpha + y_{t-2} + \varepsilon_{t} + \varepsilon_{t-1}\\
      \vdots \\
      &= y_0 + t\alpha + \sum_{j=1}^{t}\varepsilon_j  
\end{aligned}  
\end{equation}
이 때, $E(y_t) = E(y_0) + t\alpha$, $var(y_t) = t\sigma^2$가 되므로, 평균과 분산이 둘 다 $t$에 의존하게 된다.\\
\textbf{case 2) $|\beta| > 1$}\\
\begin{equation}
\begin{aligned}
y_t &= \alpha + \beta y_{t-1} + \varepsilon_t\\
    &= \alpha + \beta(\alpha + \beta y_{t-2} + \varepsilon_{t-1}) + \varepsilon_t\\ 
    &= (\alpha + \beta\alpha) + \beta^{2}y_{t-2} + \varepsilon_{t} + \beta\varepsilon_{t-1}\\
    &= \vdots\\
    &= (\alpha + \beta\alpha + \cdots + \beta^{t-1}\alpha) + \beta^{t} y_{0} + \sum_{j=0}^{t-1}\beta^{j}\varepsilon_{t-j}\\
    &= \alpha\sum_{j=0}^{t-1}\beta^{j} + \beta^{t} y_{0} + \sum_{j=0}^{t-1}\beta^{j}\varepsilon_{t-j}
\end{aligned}  
\end{equation}
AR(1) 과정 $\{y_t\}$을 식 (7)과 같이 전개하였을 때, 평균과 분산, 자기공분산은 다음과 같이 도출된다.\footnote{초기값 $y_0$은 일정한 상수로 가정한다.}\\
\begin{equation}
\begin{aligned}
E(y_{t}) &= \alpha\sum_{j=0}^{t-1}\beta^{j} + \beta^{t} y_{0}\\
         &= \frac{\alpha(1-\beta^{t})}{1-\beta} + \beta^{t} y_{0}  
\end{aligned}    
\end{equation}
여기서 $|\beta| > 1$이면 $\beta^t$의 값이 무한히 증가하므로, 평균이 시간에 따라 발산한다.
\begin{equation}
\begin{aligned}
var(y_t) &= \sum_{j=0}^{t-1}\beta^{2j}var(\varepsilon_{t-j})\\
         &= \sigma^2\sum_{j=0}^{t-1}\beta^{2j}
\end{aligned}    
\end{equation}
따라서 $var(y_t) = \sigma^2 \frac{1-\beta^{2t}}{1-\beta^2}$이다. 즉, 분산도 시간에 따라 일정하게 유지되지 못하는 것을 알 수 있다.













\end{document}
